% Section: §2 Motivated Belief at the Single-Agent Level
% PDF context: bridge section locating the paper in the empirical motivated-reasoning lineage (Kunda, Kahan, directional vs accuracy goals) and the active-inference formalisation of belief-revision cost (Hyland & Albarracin); closes by naming the single-agent --> multi-agent generalisation as the paper's open question.
% Drawing from: cluster_C1_motivated_reasoning (Kunda/Kahan lineage), cluster_C2_variational_cost (variational + bounded-cognition), and master refs.bib (hylandAlbarracin2025MindChange).
% Entries cited: 11
% Full-text papers consulted: Kahan2015_12 (Mousavi/Funder commentary surrounding Kahan's expressive-rationality piece) and Glüer-Pagin2024_11 (close philosophical reading of Kahan's identity-protective-cognition thesis).
% Key renaming TODOs:
%   Kahan2017_2     -> kahanPetersDawson2017MotivatedNumeracy   % Anchor: motivated-numeracy / PMRP design
%   Kahan2007_1     -> kahanEtAl2007WhiteMaleEffect             % Anchor: cultural cognition + white-male-effect (763 citations)
%   Kahan2017_4     -> kahan2017IdentityProtective              % Anchor: identity-protective cognition canonical statement
%   Kahan2015_12    -> kahan2015ExpressiveRationality           % Anchor: expressive-rationality framing
%   Glüer-Pagin2024_11 -> gluerSpectre2024MotivationInNumeracy  % Critical reconstruction of Kahan's mechanism
%   Bolsen2019_7    -> bolsenPalm2019MotivatedReasoning         % Field review; useful pointer to directional/accuracy distinction
%   Little2020_13   -> littleSchnakenbergTurner2021MotivatedAccountability % Formal model of directional vs accuracy
%   Savolainen2022_3 -> savolainen2022HealthMisinformation       % Motivated reasoning applied to misinformation
%   Kruglanski2020_10 -> kruglanski2020AllThinkingWishful        % Bridge from social psych to a generalised "directional cognition" claim
%   Hopkins2016_8   -> hopkins2016FreeEnergyDreaming             % Free-energy / fictive-prior gloss, used only as auxiliary AIF pointer
%   Friston2023_7   -> fristonEtAl2023FederatedInference         % Belief-sharing under FEP -- used here as the bridge into §3
% TODO refs.bib additions (foundational anchors that neither cluster surfaced):
%   kunda1990MotivatedReasoning   -- Kunda, "The case for motivated reasoning", Psych. Bull. 108(3), 1990
%   nickerson1998ConfirmationBias -- Nickerson, "Confirmation bias: a ubiquitous phenomenon", Rev. Gen. Psych. 2(2), 1998
%   lordRossLepper1979BiasedAssimilation -- Lord, Ross & Lepper, JPSP 37(11), 1979 (classic biased-assimilation experiment)

\section{Motivated Belief at the Single-Agent Level}
\label{sec:motivatedbelief}

\paragraph{Empirical motivated reasoning.}
The cognitive-psychological case that belief revision is shaped by
goals other than accuracy is long-standing and converged. The
canonical statement is Kunda's argument that reasoning is steered by
\emph{directional} as well as accuracy goals: subjects construct the
sequence of inferential steps that licenses the conclusion they would
prefer to reach, subject to the constraint that the construction
remain plausible to themselves \citep{kunda1990MotivatedReasoning}.
This framing absorbed two earlier empirical findings as special
cases: the biased-assimilation experiments of
\citet{lordRossLepper1979BiasedAssimilation}, in which the same
mixed evidence about a contested policy strengthened the prior
commitments of partisans on \emph{both} sides; and confirmation-bias
phenomena more broadly catalogued by \citet{nickerson1998ConfirmationBias}.
The empirical literature has since accumulated in two directions.
The first is Kahan's cultural-cognition programme, in which the
directional goal is identified specifically as \emph{identity
protection}: subjects process risk-relevant evidence in whatever way
preserves their standing in the affinity groups they belong to
\citep{Kahan2007_1, Kahan2017_4}. The signature empirical finding is
that the gap between cultural sub-populations \emph{widens} with
numerical and scientific sophistication rather than narrowing
\citep{Kahan2017_2}; Kahan reads this as evidence that the more
capable reasoners are more effective at the selective evaluation of
identity-incongruent evidence, a behaviour he characterises as
\emph{expressive rationality} --- rational relative to the agent's
practical stakes in group membership, irrational relative to truth
\citep{Kahan2015_12}. The mechanism remains contested:
\citet{Glüer-Pagin2024_11} argue that what Kahan's experiments
isolate is motivation \emph{to reason} rather than motivated reasoning
in the strict sense, and the broader political-science literature now
distinguishes carefully between directional and accuracy regimes and
the conditions under which one or the other dominates
\citep{Bolsen2019_7, Little2020_13, Savolainen2022_3}. For our
purposes the contested points do not matter: the robust observation
is that human belief revision exhibits a systematic, content-sensitive
asymmetry in its receptivity to evidence, and that this asymmetry is
not adequately captured by treating the agent as an unbiased Bayesian
updater.

\paragraph{Variational formalisation.}
The active-inference community has begun to give this asymmetry a
principled mathematical form. The starting point is the variational
free-energy functional, in which an agent's update from prior to
posterior incurs a complexity cost proportional to the
Kullback--Leibler divergence between the candidate posterior and the
prior \citep{Costa2020_4}. \citet{hylandAlbarracin2025MindChange}
re-foreground this complexity term as a \emph{motivational} quantity:
the cost of changing one's mind is the KL divergence from the current
belief to the candidate revision, weighted by an affective utility
that reflects how much the agent cares about retaining the prior.
Because that divergence is asymmetric, belief revision is
path-dependent: evidence that would update a neutral prior into a
particular posterior need not, played in reverse, undo a revision
already made. This is the formal counterpart of the
biased-assimilation and identity-protection findings: preferred
beliefs are not merely held more confidently, they are protected by a
mathematically explicit revision cost that contrary evidence must
overcome. \citet{Kruglanski2020_10} reaches a sympathetic conclusion
from the social-psychological side, arguing that there is no sharp
boundary between motivated and non-motivated cognition because
\emph{all} inference is directional in the sense that it is steered
by what the agent values; the variational reading makes this claim
precise by locating the directional weight in the affective utility
that multiplies the complexity term. Adjacent strands of the
active-inference literature provide the surrounding machinery on
which the multi-agent extension will draw, including federated
inference under shared world models \citep{Friston2023_7} and the
free-energy account of belief-revision dynamics more generally
\citep{Hopkins2016_8}.

\paragraph{Open question.}
The single-agent picture is now reasonably well developed, both
empirically and formally. What it does not yet capture is the
phenomenon Kuhn observed at the community scale: that scientific
sub-fields exhibit collective motivated belief, sustained across
generations, with the directional goal residing not in any individual
ego but in the institutional and social structures that bind the
community together. The active-inference treatment so far has been
single-agent --- the cost of mind-change is borne by one mind, against
its own prior. The empirical motivated-reasoning literature is in
practice multi-agent --- the identity being protected is a group
identity --- but its formalisation has remained at the level of
within-individual mechanism. The question this paper takes up is how
the variational cost of belief revision generalises to a population
of coupled inferential agents in a way that recovers what Kuhn
observed: paradigm-level persistence that is sociologically generated,
asymmetric in time, and not reducible to the directional preferences
of any single agent.
